% PP_PARAGRAPH_BANK.tex — 4 standalone Professional Practice paragraphs
% Ready to \input{} or paste into any D7 version.
% All content drawn from Apollo's verified running-context.

%% ──────────────────────────────────────────────
%% 1. TUCKMAN — Team Development Stages
%% ──────────────────────────────────────────────

% Drop into §2 (My Team) or §3 (My Impact)
Tuckman's model of group development provides a useful lens for our trajectory~\cite{tuckman1965}.
Forming was brief: the first meeting divided responsibilities clearly---Harish on hardware,
Zian on path planning, Robin on state machines, Apollo and Edward on computer vision---and
the team moved quickly into preparation.
Storming, however, was protracted and externally imposed.
Two practice flight days yielded no usable data because the provided hardware was non-functional,
while other teams flew successfully; frustration and demotivation spread across the group.
The team was blocked not by interpersonal conflict but by an inability to validate any hypothesis
without a working drone, leaving weeks of preparatory work untested.
Norming finally emerged over Easter, when hardware faults were resolved and the state machine
was confirmed on real hardware for the first time, giving the team a shared baseline of trust
in the system.
\textcolor{red}{[APOLLO: did you ever feel the team reached Performing?]}

%% ──────────────────────────────────────────────
%% 2. LENCIONI — Five Dysfunctions of a Team
%% ──────────────────────────────────────────────

% Drop into §2 (My Team) or §4 (My Support)
Lencioni's pyramid identifies \emph{absence of trust} as the foundational dysfunction from
which all others cascade~\cite{lencioni2002}.
In our team this dysfunction was not interpersonal but architectural: 799 of 807 commits
were made by a single member, meaning code ownership was never genuinely shared.
Edward and Zian completed assigned tasks competently, yet neither drove work forward
independently, which meant \emph{avoidance of accountability}---the fourth dysfunction---fell
disproportionately on Apollo and Robin as the only self-directed contributors.
Integration happened late: Robin's state machine consumed Apollo's CV output through a
defined interface built under time pressure after Easter, rather than through iterative
co-development.
The result was a system that worked on demo day, but whose bus factor of one represents
a structural vulnerability that Lencioni's model would predict from the trust deficit at
the base of the pyramid.

%% ──────────────────────────────────────────────
%% 3. RISK ASSESSMENT — Progressive Test Ladder + ALARP
%% ──────────────────────────────────────────────

% Drop into §1 (My Role) or §3 (My Impact)
The progressive test ladder---bench commands, passive detection during manual flight,
autonomous search with logging only, then full mission---was designed as a formal risk
mitigation framework, reducing both the likelihood and consequence of failure at each
stage before escalating to the next.
When the provided hardware failed on both practice flight days, this represented the
materialisation of a schedule risk that no amount of simulation could absorb: the team
entered demo day having skipped intermediate rungs of the ladder entirely.
Flight altitude decisions followed an ALARP (As Low As Reasonably Practicable) approach:
conservative parameters (\texttt{--alt 20 --conf 0.5}) were specified for first flights
to keep residual risk tolerable while still achieving detection.
The runtime geofence---which triggers emergency RTL if the drone breaches the Flight Area
polygon---was implemented as a likelihood-reduction control, ensuring that even a software
fault could not send the aircraft into the adjacent SSSI no-fly zone.

%% ──────────────────────────────────────────────
%% 4. SIMULATION AS THINKING TOOL — Graham / Innovation Tools
%% ──────────────────────────────────────────────

% Drop into §1 (My Role) or §3 (My Impact)
Graham's Innovation Tools lecture argues that tools support thinking rather than replace
it~\cite{graham_innovation}.
The full-mission simulation embodied this principle: rather than holding failure modes
in my head, I had a concrete environment to run and observe.
I ran the simulation thousands of times, deliberately injecting camera blur, GPS drift,
and edge cases such as a casualty placed near the no-fly zone boundary.
Each run externalised a question---\emph{what happens if the drone detects at the NFZ
edge?}---and returned an observable answer, converting abstract worry into tested knowledge.
Robin's suggestion of repulsive force fields for NFZ avoidance, for example, was explored
over several days of simulated experimentation rather than debated in the abstract.
When hardware failures consumed both practice flight days, the simulation became the
only venue for iterative learning; the fact that the full pipeline worked on demo day
without prior real-hardware integration owes much to the thousands of simulated rehearsals
that preceded it.
